We consider an approach from literature that is capable of learning logical rules from binary data.
It is a binary classification approach, meaning the input is a set of points,
their binary features and for each point a target boolean (True/False).
The output is a boolean formula in Disjunctive Normal Form (DNF, OR-of-ANDs).

A DNF is a disjunction (OR) of one or more conjunctions (AND).
We refer to a conjunction as simply a ``rule''.
A rule consists of one more clauses.
In the boolean setting, a clause is a literal, meaning a feature or its negation.
For instance, if $A$ is a feature representing whether TF A is active,
then the corresponding clauses are $A$ and $\neg A$ (``not A'').
If $A$ and $B$ are features, then an example of a rule is
\[(A \wedge \neg B),\]
meaning ``$A$ and not $B$''.
A disjunction of rules is what we call a ruleset.
For instance,
\[(A \wedge B) \vee (B \wedge \neg C),\]
meaning ``Either $A$ and $B$, or $B$ and not $C$'',
is a ruleset.
In practice, negated literals do not need to be modeled separately,
as a feature's negated values can be added as a separate feature during pre-processing.

The goal of the approach is to find a ruleset describing the points that have target value ``True''.
We say those points are positive points, whereas points with a target value of ``False'' are called negative points.
Negative points should not satisfy the ruleset.
Hence, the ruleset found should evaluate to ``True'' for positive points and ``False'' for negative points.
In other words, the ruleset is a binary classifier and the target values are the (binary) class labels.
The objective of this approach is to find a ruleset that maximizes accuracy, that is to maximize
\[\frac{\text{\# correctly classified points}}{\text{\text{\# points}}},\]
where ``\#'' stands for ``number of''.
Since the number of points is constant, we may multiply by it and instead maximize
\[\text{\# correctly classified points}.\]
This is equivalent to minimizing
\[\text{\# points} - \text{\# correctly classified points},\]
as we simply negate and add a constant.
This is equal to
\[\text{\# misclassified points}.\]
Separating positive and negative points, we find this equal to
\begin{equation}\label{eq:nr_of_falses}
    \text{\# false negatives} + \text{\# false positives}.
\end{equation}
Hence, maximizing accuracy and minimizing this quantity are equivalent.
In \cite{lawless_interpretable_2023}, it turns out to be easier to minimize a similar quantity, namely
\begin{equation}\label{eq:nr_of_falses_hamming}
\text{\# false negatives} + \text{\# times a rule satisfies a negative point.}
\end{equation}
For every rule in the ruleset, it counts how many negative points it covers.
Hence, we do not just count whether a negative point is misclassified, but we count how often.
It follows that \eqref{eq:nr_of_falses_hamming} is at least as large as \eqref{eq:nr_of_falses}.
If the rules do not ``overlap'' much, then \eqref{eq:nr_of_falses_hamming} is close to \eqref{eq:nr_of_falses}.
This method has been shown to work well on a variety of binary datasets.

We wish to apply this method to TF data and PD-L1 expression.
Conceptually, the target boolean should represent whether PD-L1 is ``Expressed'' (positive) or ``Not expressed'' (negative).
A simple and intuitive way to do this is via thresholding.
That is, the target considered is
\[\text{PD-L1 expression} \geq \text{PD-L1 threshold}.\]
A natural choice for the threshold is a quantile level observed in the data, i.e.\ the $0.5$ quantile, which is the median.
However, which quantile should be chosen is unclear, as there is no clear definition for a gene being expressed or not expressed.
\autoref{fig:pdl1_expression_broad} suggests any of the depicted quantiles ($0.5$, $0.6$, $0.7$, $0.8$, $0.9$)
could be a meaningful threshold.

A feature should represent whether a TF is ``Active'' or ``Inactive''
Again, these states do not have a clear definition as the data is continuous.
For features, we have the freedom to add several binary features per TF\@, for instance
\[
    \begin{split}
        \text{TF} &\geq \text{$0.1$ quantile}, \\
        \text{TF} &\geq \text{$0.2$ quantile}, \\
        \vdots & \tquad \vdots \\
        \text{TF} &\geq \text{$0.9$ quantile}. \\
    \end{split}
\]
These binary features are meant to represent different activity levels,
e.g., ``at least somewhat active'', ``at least moderately active'', et cetera.

In this chapter, we first cover the mathematical background of the method
(\autoref{sec:boolean-model-description} and \autoref{sec:boolean-pricing-problem}).
Then we briefly discuss our implementation and compare its performance to the performance of an implementation from literature
(\autoref{sec:boolean-implementation}).
Next, we discuss the methodology for a first experiment on PD-L1 and TF data (\autoref{sec:boolean-methodology})
Finally, we discuss the results (\autoref{sec:boolean-results}).
A reader who is not interested in the mathematics and the implementation may skip to \autoref{sec:boolean-methodology}.

\section{Model description}\label{sec:boolean-model-description}
We consider the supervised binary classification setting studied in~\cite{lawless_interpretable_2023}.
%We refer to this source as the ``original paper''.
What follows is a recap that closely follows this paper.
Given is a training set of $n$ points, denoted by $\I$, from an underlying (unknown) distribution.
Each point has $m$ features from set $\J$, plus a label \[(\X_i, y_i), \quad i \in \I = \{1, \dots, n\}.\]
In this setting, feature values $\X_i \in \{0,1\}^\J$ and labels
$y_i \in \{-1,1\}$ for all $i \in \I$.
A point $i$ is said to possess feature $j$ if $X_{ij} = 1$, where $X_{ij}$ is the $j$-th element of vector $\X_{i}$.
A point $i$ is said to be positive if $y_i = 1$ and negative otherwise.

Theoretically, binary features are not restrictive.
For example, categorical features can be expressed as binary features via one-hot encoding.
Real-valued features can be discretized into a number of binary features via thresholds,
i.e., binary features reflecting $X_{ij} \leq v$ or $X_{ij} \geq v$ may be introduced for various $v$.
The number of features can grow rapidly with this approach.

The objective of this model is to learn a ruleset as described earlier.
Mathematically, define the set of possible clauses as $\cL = \J$.
A rule is a subset of $\cL$.
The candidate rule set is given by  $\K = 2^\L$, i.e.\ the power set of $\L$.
Let $\K_i$ be the set of rules satisfied by point $i$, so
\begin{equation}
    \K_i = \left\{k \in \K : \underset{j \in k}{\bigwedge} X_{ij} \right\}, \quad i \in \I. \label{eq:boolean_K_i}
\end{equation}
In other words, $\K_i$ is the set of rules that cover point $i$.
Observe that $\emptyset \in \K$ and that the conjunction in \eqref{eq:boolean_K_i}
is always true for $k=\emptyset$.
So $\emptyset \in \K_i$ for any $i \in \I$.
Therefore, the ``empty rule'', $\emptyset$, covers all points under this definition.



\begin{definition}[DNF Classifier, \cite{lawless_interpretable_2023}]\label{def:dnf_classifier}
    Given a ruleset $S \subseteq \K$, the DNF classifier is defined as follows:
    \[\hat{y}_S(\textbf{X}_i) ~=~
    \begin{cases}
        1, & \text{if } \abs{\K_i \cap S} > 0, \\
        0, &\text{else.}
    \end{cases}\]
\end{definition}

There exist $2^m$ possible rules, i.e. $|\K| = 2^m$.
Therefore, as $m$ grows, enumerating all possible rules quickly becomes intractable.
In practice, one may wish to restrict $S$, for example by limiting the number of clauses to $D$.
This restricts the number of possible rules to \[\sum_{d=0}^D \binom{m}{d}.\]


The aim is to minimize the \emph{\zo classification error}, which counts the number of datapoints that are misclassified by a rule set $S$.
Hence, it is the sum of the \emph{\zo loss} of each datapoint.

\begin{definition}[\zo loss and \zo classification error]\label{def:01_loss}
    The \emph{\zo loss} of point $i$ subject to ruleset $S$ is
    \[\ell_S(\bm{X}_i) = \Ind\left(\hat{y}_S(\bm{X}_i) \neq y_i\right) =
    \begin{cases}
        \Ind(\abs{\K_i \cap S }= 0), & \text{if } y_i = 1,\\
        \Ind(\abs{\K_i \cap S} > 0), &\text{otherwise},
    \end{cases}\]
    where $\Ind(\mathcal{A})$ is the indicator function (i.e., $\Ind(\mathcal{A}) = 1$ if $\mathcal{A}$ is true and 0 otherwise).
The \emph{\zo classification error} of ruleset $S$ is
    \[\ell_S(\bm{X}) =
    \sum_{i \in \I} \ell_S(\bm{X}_i). \]
\end{definition}
Closely related is the accuracy, which is the fraction of correctly classified points.
\begin{definition}[Accuracy]
    The \emph{accuracy} of a ruleset $S$ is
    \[
        \text{Accuracy}(S) = \frac{1}{n}\left|\left\{i \in \I : \hat{y}_S(\bm{X}_i) = y_i\right\}\right|.
    \]
\end{definition}
Minimizing the \zo classification error is equivalent to maximizing the classification accuracy, as illustrated in
the introduction to this chapter.


We now formulate the problem of finding a ruleset with minimum \zo classification error as an Integer Program (IP).
Partition the points into two sets, the positive class and the negative class:
\[\cP = \{i \in \I : y_i = 1\} \text{ and } \cN = \{i \in \I : y_i \neq 1\}.\]
Let $w_k \in \{0,1\}$ be a variable indicating if rule $k \in \K$ is used in the rule set.
Let $\zeta_i\in\{0,1\}$  be a variable indicating if point $i\in \pun$ is misclassified.
The IP is given by


%\begin{alignat}{3}
%\textbf{min} && \quad \sum_{i\in \cal{P}} \zeta_i +\sum_{i\in \cal{N}} \zeta_i ~~ \label{obj:01}  \\
%	\textbf{s.t.}  && \quad
%	 \zeta_i + \sum_{k\in \K_i} w_k \geq & 1 , \quad && \forall i \in \cal{P}, \label{const:accP01} \\
%	 && w_k  \leq &   \zeta_i , \quad  && \forall i \in {\cal N}, k \in \K_i,  \label{const:accZ01} \\[.1cm]
%	%	& w_k,  ~\zeta_i  \in\{0,1\}, \quad \forall {k\in \cal{K}}, i \in {\cal P} \cup {\cal N}. ~~ \label{const:accBin01} \\
%	&& w\in \{0,1\}^{\abs{\K}},~ & \zeta\in \{0,1\}^{\abs{\cP \cup {\cal N}}} \label{const:accBin01}
%\end{alignat}

\begin{subequations}
    \label{eq:BM_orig}
    \begin{align}
        \mathbf{\min} \quad & \sum_{i \in \cP} \zeta_i + \sum_{i \in \cN} \zeta_i \label{eq:BM_orig-obj} \\
        \st \quad & \zeta_i + \sum_{k \in \K_i} w_k \geq 1, && \forall i \in \cP, \label{eq:BM_orig-misP} \\
        & w_k \leq \zeta_i, && \forall i \in \cN,\; k \in \K_i, \label{eq:BM_orig-misN} \\
        & w \in \{0, 1\}^{\K}, \quad \zeta \in \{0, 1\}^{\cP \cup \cN}. \quad && \label{eq:BM_orig-vars}
    \end{align}
\end{subequations}
The objective of this IP is exactly the \zo classification error.
To see that this IP is correct, consider first \eqref{eq:BM_orig-misP}.
Let $i \in \P$.
If $w_k = 1$ for some $k \in \K_i,$ then $i$ is covered by the ruleset.
Indeed, then $\zeta_i=0$ is valid and optimal.
If not, then $i$ is not covered and $\zeta_i = 1$ to satisfy the constraint.
In either case, $\zeta_i$ correctly represents whether $i$ is misclassified.
Now consider \eqref{eq:BM_orig-misN}.
If $w_k = 1$ for some $k \in \K_i$, then $i$ is covered and \eqref{eq:BM_orig-misN} implies $\zeta_k = 1$.
Otherwise, $\zeta_i =0$ is valid and optimal.
In either case, $\zeta_i$ correctly represents whether $i$ is misclassified.


This formulation has one integer variable per point and multiple constraints per negative datapoint.
For a large number of points, the IP becomes large.
In many practical classification tasks, the number of positive points is small compared to the total.
To limit the size of the IP, the Hamming loss may be considered.
In the Hamming loss, we do not count the number of negative points that are misclassified, but we count
the total number of times a negative point is misclassified by a rule.

\begin{definition}[Hamming loss \cite{lawless_interpretable_2023}] The \emph{Hamming loss} of point $i$ under ruleset $S$ is given by
\[\ell^{h}_S(\bm{X}_i) =
\begin{cases}
    \Ind(\abs{\K_i \cap S }= 0), & \text{if } y_i = 1,\\
    \abs{\K_i \cap S}, &\text{otherwise.}
\end{cases}\]
\end{definition}

It has been illustrated that Hamming loss is a good alternative to \zo loss in practical settings~\cite{lawless_interpretable_2023}.

Furthermore, it is useful in practice to restrict the complexity of a ruleset.
Define the complexity of a rule $k \in \K$ as
\begin{equation}
    c_k = 1 + |k|, \label{eq:boolean_ck}
\end{equation}
so one plus the number of features in the rule.
The complexity of a ruleset $S \subseteq \K$ is given by $\sum_{k\in S} c_k,$ which we bound by $C$.

The new IP that uses Hamming loss becomes

\begin{subequations}
    \label{eq:BM}
    \begin{align}
        \mathbf{\min} \quad & \sum_{i \in \cP} \zeta_i + \sum_{i \in \cN} \sum_{k \in \K_i} w_k \label{eq:BM-obj} \\
        \mathbf{\text{s.t.}} \quad & \zeta_i + \sum_{k \in \K_i} w_k \geq 1, & \forall i \in \cP, \label{eq:BM-misP} \\
        & \sum_{k \in \K} c_k w_k \leq C, & \label{eq:BM-complex} \\
        & w \in \{0, 1\}^{\K}, \zeta \in \{0, 1\}^{\cP}. \label{eq:BM-binary2}
    \end{align}
\end{subequations}
For a positive point $i$, the justification given for \eqref{eq:BM_orig} still holds.
For a negative point $i$, there is no $\zeta_i$, but instead the number of rules wrongly covering it are counted
directly in the objective.

The number of possible rules $|\K|$ grows quickly as $m$ grows.
Therefore, column generation is used to solve the LP relaxation of this problem.
A simple approach for obtaining an integer solution is employing a price-then-branch heuristic.
In this technique, the relaxed master problem (RMP) is first solved to optimality.
Next, the restricted integer master problem, with only the previously generated columns, is solved to optimality using a standard IP solver.
While it is possible to solve the problem exactly via a branch-and-price approach,
this is harder to implement, it is slower,
and it is unclear whether it improves the quality of the output significantly.
In some cases, the price-then-branch heuristic can reveal that an exact solution is
only marginally better in terms of objective.

\section{Pricing problem}\label{sec:boolean-pricing-problem}

Let $(\bm{\mu},\lambda) \in \mathbb{R}_+^{\abs{\cP}} \times \mathbb{R}_+$ be an optimal dual solution to RMP, where dual values $\bm{\mu}$, $\lambda$ correspond to constraints \eqref{eq:BM-misP} and \eqref{eq:BM-complex} respectively.
The reduced cost of a variable $w_k$, associated with a rule $k$ not yet in the master problem, can be expressed as
\begin{equation}
    \hat \rho_k =  \sum_{i \in \cN} \Ind(k \in \K_i)
    -\sum_{i \in \cP}  \mu_i  \Ind(k \in \K_i) +\lambda  c_k.  \label{eq:BP_redcost}
\end{equation}
For the derivation we refer to \autoref{sec:reduced_costs_boolean_model}.
If $\rho_k < 0$, then adding rule $k$ to the relaxed master problem can improve its objective.
If $\rho_k \geq 0$ for all $k \in \K$, we have found an optimal LP solution.
Deciding whether this is the case is called the pricing problem.

To solve the pricing problem,
we construct an IP.
Introduce the variables $z_l\in\{0,1\}, \; l\in \L$ representing whether the rule includes clause $l$.
We also introduce the variables $\delta_i\in\{0,1\}, \; i\in \I$ denoting whether point $i$ is covered by the rule.
The complexity of a rule can be expressed as
\begin{equation}\label{eq:c_z}
    c(\bm{z}) = 1 + \sum_{l \in \L} z_l, \tquad \bm{z} \in \{0,1\}^{\L},
\end{equation}
which we limit by $D$.
Define
\begin{equation}\label{eq:S_i_boolean}
    S_i = \{l \in \L \fw X_{ij} = 0\}, \tquad i \in \I.
\end{equation}
Hence, the set $S_i$ contains all clauses (binary features in this setting) that point $i$ does not satisfy (possess).
The pricing problem may be solved by minimizing \eqref{eq:BP_redcost}.
The corresponding IP, as described in~\cite{lawless_interpretable_2023}, is given by
\begin{subequations}
    \label{eq:BP}
    \begin{align}
        \min \quad & \sum_{i \in \cN} \delta_i - \sum_{i \in \cP} \mu_i \delta_i +
        \lambda c(\bm{z}) \label{eq:BP-obj} \\
        \st \quad
        & D\delta_i + \sum_{j \in S_i} z_j \leq D, && i \in \I, & \label{eq:BP-upper}\\
        & \delta_i +  \sum_{j \in S_i} z_j \geq  1, && i \in \I,  &\label{eq:BP-lower}\\
        & \sum_{j \in \J} z_j \leq D, \quad \label{eq:BP-complex}\\
        & z\in \{0,1\}^{\J},~ \delta\in \{0,1\}^{\I}.\label{eq:BP-int}\qquad
    \end{align}
\end{subequations}

If $\delta_i = 1$, then by \eqref{eq:BP-upper}, no clause that cuts off $i$ can be chosen, so $i$ must covered by the rule.
If $\delta_i = 0$, then by \eqref{eq:BP-lower}, some clause that cuts off $i$ must be chosen, so $i$ cannot be not covered by the rule.
Hence, $\delta_i$ reflects whether $i$ is covered.
Therefore, the objective \eqref{eq:BP-obj} matches the reduced cost \eqref{eq:BP_redcost}.
So the IP finds a rule with minimum reduced cost.


Constraints \eqref{eq:BP-upper} can also be expressed as
\begin{equation}\label{eq:BP-misP-better}
    \delta_i + z_j \leq 1, \quad i \in \I, \; j \in S_i.
\end{equation}
This formulation may be favorable in terms of solver efficiency, as it provides a sharper LP relaxation.
These two formulations pose a tradeoff; while \eqref{eq:BP-misP-better} implies the \eqref{eq:BP-upper},
it is also a larger family of constraints which could significantly slow down the LP solver.
Furthermore, it is unclear how the pre-solvers employed modern IP solvers handle these formulations within this IP\@.
This makes it hard to predict which formulation is favorable.

%Here, $\I^{+}$ and $\I^{-}$ may simply be replaced by $\I$. \todo{explain why / why one would use se particular sets}.

\section{Implementation}\label{sec:boolean-implementation}
A custom solver is implemented in C++ using the SCIP 10 framework with SoPlex 8~\cite{hojny_scip_2025}.
Note that the implementation of~\cite{lawless_interpretable_2023} is available on GitHub~\cite{lawless_conlawfaircg_2026}.
That implementation is writen in Python and uses proprietary software Gurobi.
Our implementation is fully open source and also available on \href{https://github.com/kalkoen/rule-sets-for-tf}{GitHub}.
This custom solver includes an efficient heuristic pricer and was designed to support the extensions discussed in the next chapters.
The custom solver has two modes:
\begin{enumerate}
    \item \textbf{Full:} Constructs and solves the full master problem \eqref{eq:BM}, with all possible rules (columns),
    \item \textbf{Column generation (CG):} Performs the price-then-branch approach.
\end{enumerate}
The Full mode can only realistically be solved for datasets with small $\J$ or small $D$.
It can serve to verify, to some extent, the performance and validity of the column generation approach.
The CG mode can use up to three pricers:
\begin{enumerate}
    \item \textbf{Iterative:} Iterates over all possible rules and computes reduced cost naively,
    \item \textbf{Exact:} Solves MIP \eqref{eq:BP} to optimality,
    \item \textbf{Heuristic:} Employs a beam search.
\end{enumerate}
The Exact pricer produces a single rule minimizing reduced cost (unless exact pricing is interrupted early).
The other pricers can be configured to return up to some maximum of rules, and to return a ``diverse'' collection of rules.

We now discuss the Heuristic pricer and the efforts taken to implement it efficiently.
The Heuristic pricer builds a tree, in which every node is a rule and each layer represents the rules of a given size.

The root node (at depth $0$) is the empty rule $\emptyset$, covering all points.
The Heuristic pricer starts by constructing all children of the root node and computing the reduced cost.
Then the nodes in this layer are pruned: only the nodes with the lowest reduced cost are kept.
The beam width at this depth of the tree determines the number of nodes kept.
This process repeats: in iteration $\nu$, the children of each node at depth $\nu-1$ are constructed.
Next, the number of nodes at level $\nu$ are reduced to the beam width of layer $\nu$.

Computing a node's reduced cost naively costs $\Theta(nd)$ operations,
where $d$ is the depth of the node, which is the rule's size.
For large datasets, this can be slow.
We note that there exists overlap between the computations of parents and those of their descendents.
We may exploit this overlap.
Let $k_\text{parent}, k_\text{child}$ be the rules corresponding to some parent and child node, respectively.
Recall that the reduced cost
\begin{equation}
        \hat \rho_{k_\text{child}} =  \sum_{i \in \cN} \Ind(k_\text{child} \in \K_i)
    -\sum_{i \in \cP}  \mu_i  \Ind(k_\text{child} \in \K_i) +\lambda  c_{k_\text{child}}.  \label{eq:BP_redcost_node}
\end{equation}
Define $\I_k$ as the set of points covered by rule $k$, so
\begin{equation}
    \I_k \coloneqq \{i \in \I \fw k \in \K_i\} = \left\{i \in \I \fw \bigwedge_{j \in k} X_{ij}\right\}, \tquad k \in \K.\label{eq:boolean_I_k}
\end{equation}
Define $\P_k$ and $ \N_k$ analogously.
Observe that
\[k_\text{parent} \notin \K_i \implies k_\text{child} \notin \K_i, \tquad i \in \I.\]
In other words, a point that does not satisfy the parent rule can never satisfy the child rule,
because the child rule is more restrictive.
Then
\[
    \I_{k_\text{child}} \subseteq \I_{k_\text{parent}}.
\]
The analogous property holds for $\P_k$ and $\N_k$.
It follows that
\begin{equation}
    \begin{split}
        \hat \rho_{k_\text{child}} = \; &  |\cN_{k_\text{child}}|
    -\sum_{i \in \cP_{k_\text{child}}}  \mu_i +\lambda c_{k_\text{child}} \\
        = \; & |\cN_{k_\text{parent}}| - |\cN_{k_\text{parent}} \cap \N_{k_\text{child}}| \\
    & -\left(\sum_{i \in \cP_{k_\text{parent}}} \mu_i - \sum_{i \in \cP_{k_\text{parent}} \cap \P_{k_\text{child}}} \mu_i \right) \\
        & +\lambda (c_{k_\text{parent}}+1) \\
        = \; & \hat{\rho}_{k_\text{parent}} - |\cN_{k_\text{parent}} \cap \N_{k_\text{child}}|
        + \sum_{i \in \cP_{k_\text{parent}}\cap \P_{k_\text{child}}} \mu_i
        + \lambda.
    \end{split}
    \label{eq:BP_redcost_node_deriv}
\end{equation}
To use this equation, we may store $\P_k$, $\N_k$ and the reduced cost for each node.
In our implementation, the sets are stored as bitmaps, reducing memory footprint and allowing for efficient treatment of
\[\cP_{k_\text{parent}} \cap \P_{k_\text{child}} \quad \text{and} \quad \N_{k_\text{parent}} \cap \N_{k_\text{child}}.\]
The cost of this approach is memory.
For every node kept in the search tree, these bitmaps need to be stored.
To avoid unnecessary memory use, these bitmaps are computed after pruning,
meaning these bitmaps are never created for pruned nodes.

The heuristic pricer can also be configured to do a ``deep'' search.
In this case, pruning is not done per layer but per parent node.
Setting the ``beam width'' (which is no longer a beam width, but a parent's degree) to $m$ or higher,
this effectively means the heuristic pricer does enumeration of all rules up to size $D$, only through this memory-heavy but computationally-efficient approach.
Depending on the available memory, the number of points number of features and $D$, the iterative pricer may be preferred.

We compiled our implementation on Eindhoven University of Technology's High Performance Cluster (HPC), allowing us to
perform many experiments simultaneously.
The HPC is a shared computer.
Hence, we cannot guarantee experiments are always run under the same load conditions.
Furthermore, we did not specifically request the same exact hardware to be used for each experiment.
Therefore, running times are not to be compared between experiments.

To test whether our implementation performs similarly to the implementation used in \cite{lawless_interpretable_2023}
we perform a benchmark.
We test our implementation on datasets marked ``small'' in the paper.
Some of the datasets could not be traced back.
The others were downloaded from the UCI Machine Learning Repository~\cite{kelly_markelle_uci_nodate}.
We set $C$ to the mean value reported in the paper, which was reportedly the average
$C$ yielding the best test accuracy.
For continuous variables, binarization is performed with $9$ quantiles.
Categorical variables are one-hot encoded.
We employ the heuristic pricer with beam widths $(100,50,50,20)$, identical to the paper.
The exact pricer is used when the heuristic pricer does not find any rule with negative reduced cost.
The pricing stage is limited to 2 hours.
The branching stage is limited to 1 hour, but in practice this limit is never reached.
Like in the original paper, 10-fold stratified cross validation is used.

\begin{adjustwidth}{-2cm}{-2cm}
\begin{table}[h]
    \begin{tabular}{l  V{1.7cm} r@{ }l r@{ }l r@{ }l r@{ }l}
\toprule
\parbox{1.7cm}{ \textbf{ Dataset}} & \textbf{$C$} & \multicolumn{2}{c}{\parbox{1.7cm}{\centering \textbf{Solve time (s)}}} & \multicolumn{2}{c}{\parbox{1.7cm}{\centering \textbf{Test acc. (\%)}}} & \multicolumn{2}{c}{\parbox{1.7cm}{\centering \textbf{Test acc. (\%) in \cite{lawless_interpretable_2023}}}} & \multicolumn{2}{c}{\parbox{1.7cm}{\centering \textbf{Train acc. (\%)}}} \\
\midrule
ILPD~\cite{ramana_ilpd_2022} & 11 & \num{6529} & (\num{1191}) & \num{69.4} & (\num{3.0}) & \num{69.6} & (\num{1.2}) & \num{78.2} & (\num{0.3}) \\
WDBC~\cite{wolberg_william_breast_1993} & 14 & \num{6859} & (\num{519}) & \num{93.3} & (\num{3.8}) & \num{94.0} & (\num{1.2}) & \num{99.0} & (\num{0.2}) \\
banknote~\cite{lohweg_banknote_2012} & 25 & \num{105} & (\num{35}) & \num{99.3} & (\num{0.7}) & \num{99.1} & (\num{0.3}) & \num{99.8} & (\num{0.1}) \\
heart~\cite{janosi_heart_1989} & 12 & \num{649} & (\num{303}) & \num{72.7} & (\num{7.7}) & \num{78.9} & (\num{2.4}) & \num{87.6} & (\num{0.9}) \\
ionosphere~\cite{sigillito_ionosphere_1989} & 13 & \num{7056} & (\num{456}) & \num{88.3} & (\num{5.1}) & \num{90.0} & (\num{1.8}) & \num{95.8} & (\num{0.6}) \\
tic-tac-toe~\cite{aha_tic-tac-toe_1991} & 32 & \num{1689} & (\num{567}) & \num{100.0} & (\num{0.0}) & \num{100.0} & (\num{0.0}) & \num{100.0} & (\num{0.0}) \\
transfusion~\cite{yeh_blood_2008} & 6 & \num{14} & (\num{14}) & \num{79.3} & (\num{3.9}) & \num{77.9} & (\num{1.4}) & \num{79.4} & (\num{0.5}) \\
\bottomrule
\end{tabular}

    \caption{Results for the benchmark between our implementation and the implementation of \cite{lawless_conlawfaircg_2026}.
    Metrics are reported in ``[Mean] ([Standard deviation])'' format. ``Accuracy'' is abbreviated to ``Acc.''.}
    \label{tab:benchmark_accuracy}
\end{table}
\end{adjustwidth}

The results are given in \autoref{tab:benchmark_accuracy}.
Test accuracies are similar for both implementations.
Not evident from the table is that our implementation is slower.
In some cases, the time limit of the pricing stage is reached.
We attribute this to the speed at which the exact pricer solves the pricing problem.

\section{Methodology}\label{sec:boolean-methodology}
As a first experiment, we use the boolean model on the full Broad dataset,
binarized in pre-processing via quantile thresholding.
The idea is to allow a moderate ruleset complexity (meaning a reasonably high $C$)
and study the performance of the ruleset in terms of accuracy.
At this point we are not interested in preventing overfitting, but merely the existence of a performant ruleset.

Let $\bm{v}$ be a vector containing measurements.
Define $Q_{\bm{v}}\colon[0,1] \to \R$ as the quantile function (or percentile function) over the measurements in $\bm{v}$.
The quantile function uses linear interpolation when given an unobserved quantile level as input value.
The target boolean for point $i$ is
\[y_i \geq Q_{\bm{y}}(Q_y),\]
where $Q_y \in [0,1]$ is the chosen target quantile level.

Note that the choice of $Q_y$ determines the balance of positive and negative samples in the dataset.
For $Q_y = 0.5,$ the target is simply whether a point's PD-L1 expression is at least the median.
By construction, this leads to a (nearly) perfectly balanced dataset.
For $Q_y = 0.9$, on the other hand, the dataset is heavily imbalanced, as only a tenth of the samples are positive.
It is important to keep this in mind when evaluating models.
Notably, the empty ruleset $\emptyset$ achieves an accuracy of $Q_y$ for all $Q_y \in [0,1]$.
This is because the empty ruleset covers no points and a fraction of $Q_y$ points is negative by construction.
Hence, a ruleset for $Q_y = 0.9$, should show at least $0.9$ accuracy, and ideally much higher.
Other performance measures that may be considered in the case of imbalanced datasets are precision and recall.
Precision measures the fraction of positively classified points that is truly positive.
Recall measures the fraction of positive points that was positively classified.
%\begin{definition}[Precision]
%    The precision of a ruleset $S$ is
%    \[
%        \text{Precision}(S) = \frac
%        {\left|\left\{i \in \I : \hat{y}_S(\bm{X}_i) = 1 \wedge y_i = 1\right\}\right|}
%        {\left|\left\{i \in \I : \hat{y}_S(\bm{X}_i) = 1\right\}\right|}
%    \]
%\end{definition}

Denote by $\nq$ the number of thresholds to use for each TF\@.
Then for some TF with observed activity score vector $\bm{a},$ we consider the binary features
\begin{equation}\label{eq:BM_TF_activity_quantiles}
    \begin{split}
        a_{\cdot} &\geq Q_{\mathbf{a}}\left(\frac{1}{n_\text{quantiles}+1}\right), \\
        a_{\cdot} &\geq Q_{\mathbf{a}}\left(\frac{2}{n_\text{quantiles}+1}\right), \\
        \vdots & \tquad \vdots \\
        a_{\cdot} &\geq Q_{\mathbf{a}}\left(\frac{n_\text{quantiles}}{n_\text{quantiles}+1}\right).
    \end{split}
\end{equation}
Here, $\cdot$ is a placeholder for a point $i$.
All these binary features combined form $\bm{X}.$

We consider two variants of the experiment.
In the first variant, only the binary features described in \eqref{eq:BM_TF_activity_quantiles} are used.
This means that rules with more than one TF represent simultaneous activity of these TF~.
The second variant also adds the reverse binary features, corresponding to the upper-bounds
\begin{equation}\label{eq:BM_TF_activity_quantiles_reverse}
    \begin{split}
        a_{\cdot} &\leq Q_{\mathbf{a}}\left(\frac{1}{n_\text{quantiles}+1}\right), \\
        a_{\cdot} &\leq Q_{\mathbf{a}}\left(\frac{2}{n_\text{quantiles}+1}\right), \\
        \vdots & \tquad \vdots \\
        a_{\cdot} &\leq Q_{\mathbf{a}}\left(\frac{n_\text{quantiles}}{n_\text{quantiles}+1}\right).
    \end{split}
\end{equation}
A rule can then be a combination of lower- and upper-bounds.
Assume for now that a TF appears at most once in a rule.
Then a rule can represent simultaneous activity of some TF and inactivity of others.
An upper-bound can be interpreted as an ``interfering'' TF that, when present,
prevents other TF from acting together.

We solve the model (via the price-then-branch heuristic) for
\[C = 20, \quad D = 3, \quad \nq = 1, 3, 5, 7, 9. \]
where $C$ is the maximum ruleset complexity and $D$ is the maximum rule size.
Recall that, in accordance with~\cite{lawless_interpretable_2023}, a rule's complexity is $1$ plus the number of features in it.
As such, a ruleset complexity of $20$ is equivalent to five rules of size three.
A ruleset complexity of $18$ is equivalent to six rules of size two.

We employ a heuristic pricer with beam widths $(1000, 500, 500)$.
We also employ the exact pricer.
We limit pricing (the finding of rules) to 23 hours and branching to 1 hour.


\section{Results}\label{sec:boolean-results}
Performance for the first variant is shown in \autoref{fig:pdl1ex1_accuracy}.
Performance barely increases as $\nq$ increases, though there appears to be some effect for lower $Q_y$.
Importantly, performance for low $Q_y$ is moderate, while performance for high $Q_y$ is bad.
For $Q_y = 0.9,$ for instance, accuracy is close to $0.9$,
while an accuracy of $0.9$ can be achieved with an empty ruleset.
The same is observed for $Q_y = 0.8$.
For $Q_y = 0.5$, however, an accuracy of over $0.70$ is achieved.

Performance for the second variant is shown in \autoref{fig:pdl1ex1neg_accuracy}.
The analysis is analogous.
Notably, the extra features yield modest increase in performance overall,
suggesting that lower-bounds may be important for finding relevant rulesets.

This method also gives results on the best possible rulesets that exist.
Its conclusion that, in a majority of the cases, there do not exist meaningfully better rulesets,
in the sense that no meaningfully higher accuracy can be reached.
For further details we refer the reader to \autoref{sec:boolean_model_appendix}.

Based on these results, choosing $Q_y = 0.5 $ or $Q_y = 0.6$ is preferred.
Nevertheless, the results are arguably underwhelming,
compared to the accuracies that are achieved on benchmark datasets.
While it may be possible to find consistent rulesets describing relevant phenomena,
rulesets of this form with moderate complexity ($C=20$) are not good predictors for thresholded PD-L1 expression.

Still, we are interested in the behavior of the model as we restrict complexity.
How much does performance drop?
What do simple rulesets look like?
Are they consistent among splits of the dataset?

To study these results, we will tailor the model to thresholds to speed up computation
and to set the scene for future extensions.


\begin{figure}[h]
    \centering
    \includesvg[width=0.8\textwidth]{gen/PDL1ex1/fig_PDL1ex1_accuracy.svg}
    \caption{Ruleset performance for the first variant of the experiment,
        only considering binary features \eqref{eq:BM_TF_activity_quantiles}.}
    \label{fig:pdl1ex1_accuracy}
\end{figure}

%\begin{figure}[h]
%    \centering
%    \includesvg[width=0.8\textwidth]{gen/PDL1ex1/fig_PDL1ex1_solve_time.svg}
%    \caption{Pricing phase solve time for the first variant of the experiment, only considering binary features \eqref{eq:BM_TF_activity_quantiles} }
%    \label{fig:pdl1ex1_solve_time}
%\end{figure}

\begin{figure}[h]
    \centering
    \includesvg[width=0.8\textwidth]{gen/PDL1ex1neg/fig_PDL1ex1neg_accuracy.svg}
    \caption{Ruleset performance for the first variant of the experiment,
        considering both features \eqref{eq:BM_TF_activity_quantiles} and \eqref{eq:BM_TF_activity_quantiles_reverse}.}
    \label{fig:pdl1ex1neg_accuracy}
\end{figure}

%\begin{figure}[h]
%    \centering
%    \includesvg[width=0.8\textwidth]{gen/PDL1ex1neg/fig_PDL1ex1neg_solve_time.svg}
%    \caption{Pricing phase solve time for the first variant of the experiment,
%        considering both features \eqref{eq:BM_TF_activity_quantiles} and \eqref{eq:BM_TF_activity_quantiles_reverse}}
%    \label{fig:pdl1ex1neg_solve_time}
%\end{figure}
